% Fixture M1 — documento realistico per i test del motore visuale.
% La forma è quella canonica del serializzatore (docToLatex): blocchi separati
% da UNA riga vuota, voci di lista rientrate di 2 spazi per livello.
\documentclass[a4paper,11pt]{article}
\usepackage[T1]{fontenc}
\usepackage[utf8]{inputenc}
\usepackage[italian]{babel}
\usepackage{geometry}
\usepackage{amsmath}
\usepackage{graphicx}
\usepackage{tikz}

% Comandi personali del documento.
\newcommand{\sigla}[1]{\textsc{#1}}
\newcommand{\datarelazione}{9 novembre 2026}

\title{Relazione sull'editor \LaTeX{} offline}
\author{Mario Rossi}
\date{\datarelazione}

\begin{document}
\maketitle

\begin{abstract}
Questa relazione descrive un editor \LaTeX{} completamente offline: il
sorgente resta l'unica verità del documento e la vista visuale è derivata.
\end{abstract}

\section{Introduzione}

\label{sec:intro}

Il documento che state leggendo serve da fixture per i test del motore
visuale: contiene tutti i costrutti del subset supportato e una sana
porzione di LaTeX grezzo che deve tornare nel sorgente byte per byte.

La formattazione inline è supportata: \textbf{grassetto}, \emph{corsivo},
\underline{sottolineato} e \texttt{monospace}, anche \textbf{annidata con
\emph{enfasi} dentro}. Il paragrafo continua su due righe % nota di revisione
senza riga vuota in mezzo: il newline resta nel testo.

I caratteri speciali si scrivono con l'escape letterale: il 100\% dei casi,
il simbolo \$ per la valuta, \# per i riferimenti, l'underscore \_ e la
e commerciale \&. La matematica inline resta un chip: $E = mc^2$ così come
\(a^2 + b^2 = c^2\), con eventuali pedici $x_i^2$.

Una nota a piè di pagina resta un chip verbatim\footnote{Le note non sono
modificabili nella vista visuale V1.}, come i riferimenti alla
Sezione~\ref{sec:liste} e le citazioni \cite{knuth1984}.

% I costrutti non supportati seguono come chip, visibili e intoccabili.

\section{Il \textbf{problema} del rendering}

Il comando personale \sigla{WYSIWYM} e il maiuscoletto \textsc{M1} non sono
nel subset: appaiono come chip inline con la slice verbatim. Lo stesso vale
per il comando \LaTeX{} e per qualunque macro non riconosciuta.

\subsection{Liste annidate}

\label{sec:liste}

Le liste sono modificabili nella vista visuale, anche annidate:

\begin{itemize}
  \item Prima voce con \textbf{grassetto}.
  \item Seconda voce con un elenco annidato:
    \begin{itemize}
      \item Sotto-voce annidata con \emph{corsivo}.
      \item Altra sotto-voce con matematica $x^2$ come chip.
    \end{itemize}
  \item Terza voce con passi numerati:
    \begin{enumerate}
      \item Aprire il progetto.
      \item Premere F5 per compilare.
    \end{enumerate}
  \item Quarta e ultima voce.
\end{itemize}

Anche l'elenco numerato di primo livello è un blocco modificabile:

\begin{enumerate}
  \item Compilare il documento con F5.
  \item Verificare il PDF generato.
  \item Condividere il sorgente con i collaboratori.
\end{enumerate}

\subsection{Citazioni e codice}

Le citazioni sono ambienti supportati con più paragrafi:

\begin{quote}
La semplicità è la massima raffinatezza.

Il codice è poesia che il computer può eseguire, con \emph{enfasi} dove serve.
\end{quote}

Il blocco verbatim è modificabile e arriva tale e quale al sorgente:

\begin{verbatim}
$ pdflatex main.tex
This is pdfTeX (ribadito: <nessun> download & 100% offline)
\end{verbatim}

\section*{Nota del curatore}

Le sezioni con la forma stellata sono supportate dal subset: la stellina
viene conservata nel round-trip e mostrata come attributo del titolo.

\section{Matematica e blocchi grezzi}

\label{sec:matematica}

La matematica display non è resa in V1: resta un chip blocco verbatim, come
la matematica inline. La TeXcia ride.

\[
V = \frac{4}{3}\pi r^3
\]

\begin{equation}
E = mc^2 \label{eq:emc}
\end{equation}

\begin{align}
a &= b + c \\
d &= e + f \nonumber
\end{align}

\subsection{Figure e tabelle}

Figure, tabelle e disegni non sono nel subset: interi ambienti come chip.

\begin{figure}[htbp]
  \centering
  \includegraphics[width=0.7\textwidth]{grafico.pdf}
  \caption{Un grafico di esempio con label.}
  \label{fig:grafico}
\end{figure}

\begin{table}[htbp]
  \centering
  \begin{tabular}{lr}
    Simbolo & Valore \\
    \hline
    $\alpha$ & 1.618 \\
    $\pi$ & 3.14159 \\
  \end{tabular}
  \caption{Costanti usate negli esempi.}
\end{table}

\begin{center}
Testo centrato, con \textbf{enfasi} che resta grezza nel chip.
\end{center}

\begin{tikzpicture}[scale=1.0]
  \draw[thick] (0,0) rectangle (3,2);
  \node at (1.5,1) {Editor};
\end{tikzpicture}

\section{Compatibilità e limiti del subset}

Il subset V1 copre il testo di uso quotidiano: paragrafi, titoli, liste,
citazioni, verbatim e la formattazione inline. Tutto il resto — matematica,
figure, tabelle, disegni — viaggia come chip grezzo che il round-trip
riporta identico a se stesso.

I limiti sono dichiarati, non nascosti: la vista visuale non inventa mai
LaTeX che l'utente non ha scritto e non scarta mai LaTeX che non capisce.
Un documento con moltissima matematica resta principalmente un documento da
editare in Sorgente, e va benissimo così.

\begin{itemize}
  \item La vista visuale deriva dal sorgente: mai due modelli sincronizzati.
  \item Il toggle è innocuo quando il serialize riproduce il sorgente byte
    per byte, come in questo file.
  \item La matematica renderizzata arriva con un milestone successivo.
\end{itemize}

Un ultimo esempio di costrutti ai margini: l'elenco descrittivo è un chip,

\begin{description}
  \item[WYSIWYM] What You See Is What You Mean.
  \item[Chip] porzione di sorgente mostrata grezza e non modificabile.
\end{description}

mentre la nota al piede resta tale\footnote{Anche qui, slice verbatim.} e il
simbolo tilde nell'indirizzo utente~2 è testo valido.

\section{Conclusioni}

Il round-trip di questo documento deve essere un'identità byte per byte:
paragrafi, liste, citazioni, verbatim, titoli e chip tornano al loro posto.
La coda dopo il comando di chiusura contiene un commento finale.
\end{document}
% Fine del file: fixture generata per i test M1.
